% Standalone chunk C: Sub-Gaussian probes via truncation
% (POST-FIX, Round 2).
% Covers: Prop truncation (Gaussian tail event),
%         Prop theta_bound (truncation-bias offset bound),
%         formalized with explicit moment-bound derivation.
% Source: conf.tex lines ~2157-2258 (post-fix).
% Shared preamble for standalone chunks. Do not compile directly; \input{} it.
\documentclass{article}
\usepackage[utf8]{inputenc}
\usepackage[T1]{fontenc}
\usepackage{lmodern}
\usepackage[margin=1in]{geometry}
\usepackage{amsmath,amssymb,amsfonts,amsthm,bm,mathtools}
\usepackage{enumitem}
\usepackage{microtype}
\usepackage{xcolor}
\usepackage{natbib}
\usepackage{booktabs}
\usepackage{hyperref}

\newcommand{\R}{\mathbb{R}}
\newcommand{\N}{\mathbb{N}}
\newcommand{\E}{\mathbb{E}}
\newcommand{\Prob}{\mathbb{P}}
\newcommand{\op}{\mathrm{op}}
\DeclareMathOperator{\sign}{sign}
\newcommand{\cA}{\mathcal{A}}
\newcommand{\cF}{\mathcal{F}}
\newcommand{\cH}{\mathcal{H}}
\newcommand{\cT}{\mathcal{T}}
\newcommand{\cI}{\mathcal{I}}
\newcommand{\cW}{\mathcal{W}}
\newcommand{\cTprobe}{\mathcal{T}_{\mathrm{probe}}}
\newcommand{\cE}{\mathcal{E}}
\newcommand{\norm}[1]{\left\lVert #1 \right\rVert}
\newcommand{\tilO}{\widetilde{\mathcal{O}}}
\newcommand{\DynReg}{\mathrm{DynReg}}
\newcommand{\rank}{\mathrm{rank}}
\newcommand{\range}{\mathrm{range}}
\DeclareMathOperator{\TopEig}{TopEig}
\DeclareMathOperator*{\argmax}{arg\,max}
\DeclareMathOperator{\tr}{tr}

\newtheorem{assumption}{Assumption}
\newtheorem{theorem}{Theorem}
\newtheorem{corollary}{Corollary}
\newtheorem{proposition}{Proposition}
\newtheorem{lemma}{Lemma}
\newtheorem{remark}{Remark}


\title{Chunk C: Sub-Gaussian probes via truncation\\(post-fix, Round 2 verification)}
\author{(excerpt for standalone review)}
\date{}

\begin{document}
\maketitle

\textbf{Setup context.}
\begin{itemize}[leftmargin=*,itemsep=1pt,topsep=1pt]
\item The algorithm computes $G_t=\mathcal K^{-1}(s_t u_tu_t^\top)$
using the untruncated Gaussian closed-form $\mathcal K^{-1}$.
\item Matrix Freedman requires an a.s.\ bound on increments; this only
holds on the observable event $\cA_t = \{\|u_t\|\le L\} \cap \{|y_t|\le L_y\}$
(Lem.\ G\_bound\_conf).
\item Thm.\ matrix\_bernstein\_conf uses a \emph{predictably-truncated}
MDS $\tilde X_t := \tilde G_t - \mathbb E[\tilde G_t|\cH_{t-1}]$ with
$\tilde G_t = G_t\mathbf 1\{\cA_t\}$, and leaves a deterministic bias
offset $\Theta_k := m_k^{-1}\sum_{t\in\cT_k}\E[G_t\mathbf 1\{\cA_t^c\}|\cH_{t-1}]$.
This chunk bounds $\|\Theta_k\|_\op$.
\end{itemize}

\textbf{What changed since Round~1 (delta highlights).}
\begin{itemize}[leftmargin=*,itemsep=1pt,topsep=1pt]
\item \emph{Moment bound formalized (Bug~7 fix).} Old proof used the
hand-wavy $C_s=O(S_w^4+\sigma_\varepsilon^4+\hat\sigma^4)$; fixed with
explicit binomial expansion $(a+b+c)^4\le 27(a^4+b^4+c^4)$, Gaussian
moment $\E[(u^\top\theta)^8]=105\|\theta\|^8$, and sub-Gaussian
$\E[\varepsilon_t^8]\le\text{const}\cdot\sigma_\varepsilon^8$.
\item \emph{Proposition environment.} Promoted from paragraph-level
claim to a proper Prop.~\ref{prop:theta_bound} with full proof, so
that it can be referenced by Cor.\ projector\_conf.
\end{itemize}

\section*{Proposition (Gaussian truncation)}

\begin{proposition}[Truncation of Gaussian probes]
\label{prop:truncation}
Let $u_t\sim\mathcal N(0,I_d)$ and set $L := \sqrt{2d\log(T/\delta)}$.
Define the truncation event $\cE := \bigcap_{t\in\cTprobe}\{\|u_t\|\le L\}$.
Then $\Pr(\cE^c)\le\delta$, and on $\cE$ the analysis of chunk~A.2
applies verbatim with this $L$. All subsequent results hold with an
additive $\delta$ added to the failure probability.
\end{proposition}
\begin{proof}
For $u\sim\mathcal N(0,I_d)$, $\|u\|^2$ is $\chi^2_d$; Laurent--Massart
tail $\Pr(\|u\|^2 > d + 2\sqrt{dx} + 2x)\le e^{-x}$, so choosing
$x = \log(T/\delta)$ yields $\Pr(\|u\|>\sqrt{2d\log(T/\delta)}+\sqrt d)\le\delta/T$.
Absorbing the $\sqrt d$ offset into $L$ and union-bounding over $|\cTprobe|\le T$
gives $\Pr(\cE^c)\le\delta$. On $\cE$, $\|u_t\|\le L$ a.s., which is
what Lem.\ G\_bound\_conf and Thm.\ matrix\_bernstein\_conf use. The
only $L$-dependence in the regret bound is through $R_X$ and
$C_{\mathrm{sub}}$, which pick up a $\log(T/\delta)$ factor absorbed
into $\tilO(\cdot)$.
\end{proof}

\section*{Bound on $\|\Theta_k\|_\op$}

\begin{proposition}[Truncation-bias bound]
\label{prop:theta_bound}
The deterministic bias offset
$\Theta_k := m_k^{-1}\sum_{t\in\cT_k}\E[G_t\mathbf 1\{\cA_t^c\}\mid\cH_{t-1}]$
satisfies $\|\Theta_k\|_\op \le O(d\sqrt{\delta/T})$, which is $o(1)$
in the regime of interest ($\delta$ polynomially small in $T$) and
dominated by the Freedman bound $O(R_X\sqrt{\log(2d/\delta)/m_k})$
for $m_k = O(T^{2/3})$ (the allocation used in Thm.~1).
\end{proposition}
\begin{proof}
\emph{Per-round bound.} Write $\cA_t^c \subseteq \{\|u_t\|>L\}\cup\{|y_t|>L_y\}$,
and recall (Lem.\ G\_bound\_conf's proof) that on $\{\|u_t\|\le L\}$,
$\{|y_t|>L_y\}\subseteq\{|\varepsilon_t|>L_\varepsilon\}$ by triangle
inequality. By Cauchy--Schwarz in the operator norm,
\[
  \|\E[G_t\mathbf 1\{\cA_t^c\}\mid\cH_{t-1}]\|_\op
  \;\le\; \E[\|G_t\|_\op^2\mid\cH_{t-1}]^{1/2}\cdot\Pr(\cA_t^c\mid\cH_{t-1})^{1/2}.
\]

\emph{Second-moment term.} $\|G_t\|_\op\le|s_t|\,\|u_t\|_2^2$ (using
$\|\mathcal K^{-1}\|_\op\le 1$ and $\|u_tu_t^\top\|_\op=\|u_t\|^2$).
Also
$|s_t|\le y_t^2+\hat\sigma^2\le 2(u_t^\top\theta_t)^2+2\varepsilon_t^2+\hat\sigma^2$
(using $y_t^2\le 2(u_t^\top\theta_t)^2+2\varepsilon_t^2$ by AM--GM and
$s_t=y_t^2-\hat\sigma^2$, $|s_t|\le y_t^2+\hat\sigma^2$).
We compute $\E[s_t^4\mid\cH_{t-1}]$ explicitly. Since
$u_t\sim\mathcal N(0,I_d)$ and $\|\theta_t\|\le S_w$,
$u_t^\top\theta_t\sim\mathcal N(0,\|\theta_t\|^2)$, so
$\E[(u_t^\top\theta_t)^8]=105\|\theta_t\|^8\le 105 S_w^8$ (Gaussian
moment, constant in $d$). Similarly
$\E[\varepsilon_t^8]\le\text{const}\cdot\sigma_\varepsilon^8$ by
sub-Gaussian $\psi_2$-norm. Expanding $s_t^4$ with
$(a+b+c)^4\le 27(a^4+b^4+c^4)$:
\[
\E[s_t^4\mid\cH_{t-1}]\le 27\bigl(16\E[(u_t^\top\theta_t)^8]+16\E[\varepsilon_t^8]+\hat\sigma^8\bigr)
\le C\bigl(S_w^8+\sigma_\varepsilon^8+\hat\sigma^8\bigr)=:C_s^\prime,
\]
a constant in $d$. And $\E[\|u_t\|^8]=d(d+2)(d+4)(d+6)\le C\,d^4$. By
Cauchy--Schwarz,
\[
\E[\|G_t\|_\op^2\mid\cH_{t-1}]\le\E[s_t^2\|u_t\|^4\mid\cH_{t-1}]
\le\sqrt{\E[s_t^4]\cdot\E[\|u_t\|^8]}\le\sqrt{C_s^\prime\cdot C d^4}=O(d^2).
\]

\emph{Probability term.}
$\Pr(\cA_t^c\mid\cH_{t-1})\le\Pr(\|u_t\|>L)+\Pr(|\varepsilon_t|>L_\varepsilon)\le 2\delta/T$
(by Lem.\ G\_bound\_conf).

\emph{Combining.}
\[
  \|\E[G_t\mathbf 1\{\cA_t^c\}\mid\cH_{t-1}]\|_\op
  \;\le\; O(d)\cdot\sqrt{2\delta/T}
  \;=\; O(d\sqrt{\delta/T}).
\]
Averaging over $t\in\cT_k$: $\|\Theta_k\|_\op\le O(d\sqrt{\delta/T})$,
which is $o(1)$ for any regime with $\delta$ polynomial in $1/T$.

\emph{Domination.} Compare to the Freedman bound
$O(R_X\sqrt{\log(2d/\delta)/m_k})$ with $R_X=L^2 R_s+S_w^2=O(d\log^2 T)$:
$\|\Theta_k\|_\op\ll$ Freedman iff $d\sqrt{\delta/T}\ll d\log^{5/2}T/\sqrt{m_k}$,
i.e.\ $m_k\ll T\log^5/\delta$. For $\delta=1/\mathrm{poly}(T)$ and
$m_k=O(T^{2/3})$, this is trivially satisfied.
\end{proof}

\end{document}
